\DTMsavedate{mydate}{2022-06-30}
\maketitle{Software Engineering}{Electronic Engineering}{\DTMusedate{mydate}}

% ----------------------------------------------------- %
% COURSE INFO
\subsection*{Course Information}
\vspace{0ex}%
\noindent\parbox{0.5\textwidth}{%
    \noindent\begin{tabular}{@{}ll}
                 \textsf{Course:}          & 	Software Engineering \\
                 \textsf{Course Code:}         & 	TEE 3132    \\
                 \textsf{Credits:}    & 	10
    \end{tabular}}
\vspace{2ex}\hrule\vspace{2ex}

% ----------------------------------------------------- %
% INSTRUCTOR INFO
\subsection*{Lecturer Information}
\noindent\parbox{0.5\textwidth}{%
    \noindent\begin{tabular}{@{}ll}
                 \textsf{Name:}         &    Fidelis Nhenga Mugarisanwa \\
% \textsf{Phone:}        &    \phone          \\
\textsf{Email:}        &    fidelis.mugarisanwa@nust.ac.zw \\
\textsf{Qualifications:}        & BSc in Computer Science
    \end{tabular}}
\vspace{2ex}\hrule\vspace{2ex}


% ----------------------------------------------------- %
% LEARNING OBJECTIVES
\subsection*{Learning Objectives}
The following course objectives provide a general overview of outcome expectations which will depend on the student’s motivation and work rate. Thus the following list is just an eye opener.
The course will be expected to:
\begin{outline}
    \1 Introduce the students to the Java programming language, its syntax, and file organisation.

    \1 Introduce students to Java’s core libraries.

    \1 Promote  the use of the object oriented programming paradigm among the students.

    \1 Demonstrate the potential of  the use of Java as a ubiquitous effective problem solving tool.

    \1 Engage students with real world examples.

    \1 Teach students software engineering design concepts.

    \1 Give students practical guidance to:
    \begin{itemize}
        \item coding and modifying existing code
        \item good and bad programming techniques
    \end{itemize}
\end{outline}
\vspace{2ex}\hrule\vspace{2ex}

% ----------------------------------------------------- %
% COURSE TOPICS
\subsection*{Topics}
\begin{outline}
    \1 Introduction
    \1 Java Basics Overview
    \1 Using Classes and Objects
    \1 Programming Flow Control Using Decisions
    \1 Iterations
    \1 Methods and Classes
    \1 Arrays
    \1 Inheritance and Polymorphism
    \1 Java Inbuilt Class Examples
    \1 Exceptions
    \1 Threads
\end{outline}
\vspace{2ex} \hrule\vspace{2ex}

% ----------------------------------------------------- %
% Students Assenssment
\subsection*{Assessment of Students}
\begin{outline}
    \1 Continuous assessment will consist of two written assignments and two written in-class tests.
    \1 The final assessment shall be a written closed book examination at the end of the semester.
\end{outline}
\vspace{2ex}\hrule\vspace{2ex}

% ----------------------------------------------------- %
% Grade Evaluation
\subsection*{Course Evaluation and Grading Policies}
Grades are based on:
Continuous assessment (\qty{25}{\percent}),
Final Exam (\qty{75}{\percent})
\vspace{2ex}\hrule\vspace{2ex}
% ----------------------------------------------------- %
% EQUIPMENT
% https://sites.udel.edu/computing-purchases/personal-specs/
% \subsection*{Essential Equipment}

% \vspace{2ex}\hrule\vspace{2ex}

% Policies and Other information
\subsection*{Policies and Other Information}
% Key Text or some Books
\subsection*{Key Text}
\begin{itemize}
    \item Course notes and other materials which may either be supplied or referred to during the course
    \item Schildt, H. (2015). Java The complete Reference 10th Ed Oracle Press McGraw Hill Education
    \item Cahoon J. P., & Davidson, J. W. (2004). Java 1.5 Design, McGraw Hill Higher Education
    \item Van der Linden, P. (1996). Just Java SunSoft Press Prentice Hall
    \item \href{https://www.codejava.net/}{https://www.codejava.net/}
    \item \href{https://www.edureka.co/}{https://www.edureka.co/}
    \item \href{https://www.softwaretestinghelp.com/java/}{https://www.softwaretestinghelp.com/java/}
    \item \href{http://tutorials.jenkov.com/java/}{http://tutorials.jenkov.com/java/}
\end{itemize}
\vspace{2ex}\hrule
\vspace{2ex}\hrule\vspace{2ex}